\documentclass[12pt]{article}
\usepackage[utf8]{inputenc}
\usepackage[T1]{fontenc}
\usepackage{hyperref}
\title{Vatrata Language Resources for K-12 Curriculum Development}
\author{Compiled by Claude Sonnet 5 (claude-sonnet-5) for Andr\'as Kornai}
\date{}
\begin{document}
\maketitle

\section{Introduction}
Vatrata (ISO 639-3 vlr, Glottocode vatr1237) is a retired code, and the language it named cannot be taught from the code alone. According to Glottolog it was split into Vera'a (vra) and Lemerig (lrz). Both are spoken in the Banks Islands of northern Vanuatu. The Glottolog change-request excerpt says that ``Vatrata'' was a name chosen by 19th-century missionaries on the basis of Mota and was never used by the local community. Vera'a is spoken on the west coast of Vanua Lava. Wikipedia gives about 500 speakers, and UNESCO classes it as definitely endangered. Lemerig is reported to have very few speakers (the excerpt says 2). The family label ``Bookkeeping'' in the prompt does not describe a real language family and appears to be a placeholder for a code that is no longer in use. I found no resources specific to ``Vatrata'' as a separate language.

\section{Primary Educational Resources}
I found no primers, textbooks or school programmes for Vatrata, Vera'a or Lemerig.

\section{Linguistic and Reference Materials}
Resources for Vera'a (vra), which the sources say is also known as Vatrata:
\begin{itemize}
\item Vera'a language, Wikipedia (mentions Stefan Schnell's 2011 description): \url{https://en.wikipedia.org/wiki/Vera%CA%BCa_language}
\item Glottolog entry for the retired code vatr1237: \url{https://glottolog.org/resource/languoid/id/vatr1237}
\item Austronesian Basic Vocabulary Database, Vera'a (Vatrata): \url{https://abvd.eva.mpg.de/austronesian/language.php?id=1077}
\end{itemize}

\section{Religious and Community Resources}
\begin{itemize}
\item Joshua Project, Vatrata people group, Vanua Lava: \url{https://joshuaproject.net/people_groups/15777}
\item Joshua Project, language page for vra: \url{https://www.joshuaproject.net/languages/vra}
\end{itemize}

\section{Digital and Audio Resources}
\begin{itemize}
\item Global Recordings Network, Vera'a language page: \url{https://globalrecordings.net/en/language/vra}
\end{itemize}

\section{Curriculum Development Context}
The identification must be settled first: any work should target Vera'a (vra) and/or Lemerig (lrz) with community input. The searches found no school curriculum for either language. For Vera'a the documentation includes a published grammar by Schnell (not located as a full text here) and word lists. For Lemerig the sources point to a language close to dying out.

\end{document}